% !TEX root = chapter-standalone.tex

\section{Strictification}
\label{sec:strictification}

\todotext{write this up}

\subsection{Motivation: strict versus non-strict}

In a \SY{strict monoidal category}~$\tup{\CatC, \mtimescat, \idmoncat}$, one has equalities such as
\begin{equation}
    (\Obja \mtimescatob \Objb) \mtimescatob \Objc = \Obja \mtimescatob (\Objb \mtimescatob \Objc) \qqand \Obja \mtimescatob \idmoncat = \Obja = \idmoncat \mtimescatob \Obja
\end{equation}
for all objects~$\Obja, \Objb, \Objc$.
This is very convenient: it means we never have to worry about how we parenthesize a monoidal product of several objects, nor about inserting or removing monoidal units.

In a general (non-strict) \SY{monoidal category}, these equalities are replaced by \SY{natural isomorphisms}: the \SY{associator}~$\associator_{\Obja,\Objb,\Objc}$, the \SY{left unitor}~$\leftunitor_\Obja$, and the \SY{right unitor}~$\rightunitor_\Obja$.
These are not equalities but rather morphisms that we must compose with whenever we want to ``re-parenthesize'' or ``insert/remove the unit''.
While the \SY{pentagon identities} and \SY{triangle identities} ensure that everything is coherent (roughly: any two ways of re-parenthesizing give the same result), keeping track of all the associators and unitors can be quite cumbersome.

This raises a natural question: can we simply ``pretend'' that every \SY{monoidal category} is strict, and not worry about associators and unitors at all?

\subsection{Mac Lane's strictification theorem}

The answer, remarkably, is essentially yes.
The key result is due to Mac Lane, and it says that every \SY{monoidal category} is ``equivalent'' to a \SY{strict monoidal category}, in a sense that respects the monoidal structure.
Let us make this precise.

We first need to say what it means for an \SY{equivalence of categories} to be compatible with the monoidal structures on both sides.

\begin{ctdefinition}[Monoidal equivalence]
    \label{def:monoidal-equivalence}
    Let~$\tupp{\CatC,\mtimesC,\idmoncat_{\CatC}}$ and~$\tupp{\CatD,{{\mtimesD}},\idmoncat_{\CatD}}$ be \SY{monoidal categories}.
    A \maindef{monoidal equivalence} between them is an \SY{equivalence of categories} between \CatC and \CatD in which both functors~$\funl \colon \CatC \fto \CatD$ and~$\funr \colon \CatD \fto \CatC$ are \SY{strong monoidal functors}.

    If a \SY{monoidal equivalence} exists between \CatC and \CatD, we say that they are \emph{monoidally equivalent}.
\end{ctdefinition}

The idea is that monoidally equivalent categories have ``the same'' monoidal structure, even though the categories themselves may look quite different.
In particular, any categorical construction or property that can be expressed purely in terms of the monoidal structure will transfer from one to the other.

\begin{theorem}[Mac Lane's strictification theorem]
    \label{thm:maclane-strictification}
    For every \SY{monoidal category}~$\tup{\CatC, \mtimescat, \idmoncat, \associator, \leftunitor, \rightunitor}$, there exists a \SY{strict monoidal category}~$\tup{\CatC', \mtimescat', \idmoncat'}$ and a \SY{monoidal equivalence} between~\CatC and~$\CatC'$.
\end{theorem}

The strict monoidal category~$\CatC'$ is called a \emph{strictification} of~\CatC.
The proof, which we omit here, is constructive: given~\CatC, one can explicitly build~$\CatC'$ along with the monoidal equivalence.
We will see concrete examples of such constructions in the next section.

\begin{remark}
    Mac Lane's strictification theorem is the formal justification for a common practice in the literature: when working with monoidal categories, one often suppresses the associators and unitors from the notation and writes expressions like~$\Obja \mtimescatob \Objb \mtimescatob \Objc$ without specifying the parenthesization.
    The theorem guarantees that this informal practice does not lead to errors, because any result proved in the strict setting transfers, via the monoidal equivalence, to the non-strict one.
\end{remark}

\begin{remark}
    The same result holds in the symmetric case: every \SY{symmetric monoidal category} is monoidally equivalent to a \emph{symmetric} \SY{strict monoidal category}, where the strictification also preserves the symmetry.
    However, note that strictification only makes the associator and unitors into equalities.
    The braiding~$\braiding_{\Obja,\Objb} \colon \Obja \mtimescatob \Objb \mtoiso \Objb \mtimescatob \Obja$ in general remains a non-trivial isomorphism even in the strictified version (unless~$\braiding$ is itself the identity, which would mean that the monoidal product is commutative on the nose).
\end{remark}





\subsection{Some practical strictifications}


\subsection{The category \SetL}
\label{sec:cartcatset}

We will define a category~\SetL whose objects are ``sets of tuples''.

%\subsection{Lists}
%
%Before defining L-sets, let's first recall our notion of lists.
%Given three things, for example (call them $\ela$, $\elb$, and $\elc$) we can list them in some order:
%\begin{equation}
%    [\elb, \elc, \ela].
%\end{equation}
%The square brackets are our notation for lists, and recall that the ordering matters for lists (for instance, $[\elc, \ela, \elb]$ is a different list of those same three things).
%
%Any list has a length; for example $[\elb, \elc, \ela]$ has length 3.
%For us, by definition, all lists have finite length.
%We also define such a thing as an \emph{empty list}, which we denote by $[ \ ]$.
%It has length zero.
%
%The things in a list -- for example, $\elb$, $\elc$, and $\ela$ -- are the \emph{entries} of the list.
%They each come with an index corresponding to their placement in the list (so we can speak of the 1st entry, 2nd entry, and so on).

Given sets~$\setAn{1}, \setAn{2}, \ldots, \setAn{n}$, where $n \setin \natnumbers$, we define
\begin{equation}
    \label{eq:sets-of-tuples}
    \cObj{\setAn{1}, \setAn{2}, \ldots, \setAn{n}} \definedas \makeset{ \tup{\elna{1}, \elna{2}, \ldots, \elna{n}} \mid \elna{1} \setin \setAn{1}, \elna{2} \setin \setAn{2}, \ldots, \elna{n} \setin \setAn{n} }.
\end{equation}
So for example,
\begin{equation}
    \label{eq:setL-example}
    \cObj{\wnumbers, \natnumbers, \reals} = \makeset{ \tup{\ela, \elb, \elc} \mid \ela \setin \wnumbers, \elb \setin \natnumbers, \elc \setin \reals }.
\end{equation}

% A special case is when any of the sets $\setAn{1}, \dots, \setAn{n}$ are the empty set $\Emptyset$; then we have
% \begin{equation}
%     \cObj{\setAn{1}, \setAn{2}, \ldots, \setAn{n}} = \Emptyset.
% \end{equation}
Note the case $n = 0$, where
\begin{equation}
    \cObj{} = \makeset{ \tup{} }
\end{equation}
is a singleton set whose element is the empty tuple.

% We make the convention that~$\cObj{\setAn{1}, \setAn{2}, \ldots, \setAn{n}}$ denotes the empty set if any of the sets~$\setAn{1}, \setAn{2}, \ldots , \setAn{n}$ are equal to the empty set.
% In particular, we have for example~$\Emptyset = \cObj{\Emptyset}$.

All the objects of~$\SetL$ are in particular just sets, so we can package them into a category whose morphisms are simply functions between such sets.

\begin{ctdefinition}[\SetL]\label{def:SetL}
    \SYNDEF{category of tuple-sets and functions}
    The category \SetL is the category whose objects are those sets which are of the form~$\cObj{\setAn{1}, \setAn{2}, \ldots, \setAn{n}}$, for any~$n \setin \natnumbers$ and any sets~$\setAn{1}, \setAn{2}, \ldots, \setAn{n}$.
    Morphisms in \SetL are any functions between such sets, and the composition operations and identities are the usual ones for functions.
\end{ctdefinition}

\paragraph{From \Set to~\SetL}

Observe that we can turn any set into an object of \SetL: given a set~\setA, there is the associated set
\begin{equation}
    \label{eq:setintosetl}
    \cObj{\setA} = \makeset{ \tup{\ela} \mid \ela \setin \setA }
\end{equation}
whose elements are precisely all the tuples of length one whose entry is an element of~\setA.
This defines a function
\begin{equation}
    \label{eq:setintosetlfunction}
    \begin{aligned}
        \Obof{\Set} & \to \Obof{\SetL} \\
        \setA       & \mapsto \cObj{\setA}
    \end{aligned}
\end{equation}
%\Ob_{\Set} \to \Ob_{\SetL}, \setA \mapsto \cObj{\setA}
which one might call ``bracket''.

Similarly, given a function~$\mapa \colon \setA \sto \setB$ between sets, there is an associated function~$\cMor{ \mapa }\colon \cObj{\setA} \mto \cObj{\setB}$ given by
\begin{equation}
    \label{eq:functionintosetlmor}
    \cMor{ \mapa } ( \tup{\ela }) \definedas \tup{\mapa(\ela)}.
\end{equation}
This defines a function
\begin{equation}
    \label{eq:functionintosetlmorfun}
    \Hom_{\Set}(\setA,\setB) \to \Hom_{\SetL}(\cObj{\setA},\cObj{\setB})
\end{equation}
for any sets~\setA and~\setB; one might also call it ``bracket''.

\paragraph{Concatenation of tuples}

Given tuples~$\tup{\ela, \elb, \elc}$ and~$\tup{\eld, \ele}$ we can stick them together to make the longer tuple~$\tup{\ela, \elb, \elc, \eld, \ele}$, using the already defined \SY{concatenation of tuples}.

%We call this operation \emph{concatenation}, and in the context of \SetL we think of it as a kind of multiplication.
%Our notation for concatenation of tuples is
%\begin{equation}\label{eq:concatenation-of-tuples-SetL}
%    \tup{\ela, \elb, \elc} \tupconcat \tup{\eld, \ele} = \tup{\ela, \elb, \elc, \eld, \ele}.
%\end{equation}
%\todotext{\alphubel: We already defined the concatenation of tuples...}
%For the empty tuple~$\tup{ }$ we set
%\begin{equation}\label{eq:empty-tuple-concat-1}
%    \tup{} \tupconcat \tup{\ela, \elb, \elc}  = \tup{\ela, \elb, \elc}
%\end{equation}
%and
%\begin{equation}\label{eq:empty-tuple-concat-2}
%    \tup{\ela, \elb, \elc}  \tupconcat \tup{} = \tup{\ela, \elb, \elc}.
%\end{equation}

\paragraph{Multiplication in \SetL}
\label{sec:multiplication-for-SetL}
We define a multiplication for objects of \SetL.
The symbol~$\cprod$ will denote this multiplication in infix notation.

Given~$\cObj{\setAn{1}, \ldots, \setAn{m}}$ and~$\cObj{\setBn{1}, \ldots, \setBn{n}}$ we define the operation
\begin{equation}
    \label{eq:mult-setl-sign}
    \cprod \colon \Obof{\SetL} \cartprod \Obof{\SetL} \to \Obof{\SetL}
\end{equation}
by
\begin{equation}
    \label{eq:strict-monoidal-product-sets}
    \makecprod{
        \cObj{\setAn{1}, \ldots, \setAn{m}},
        \cObj{\setBn{1}, \ldots, \setBn{n}}
    } \definedas \cObj{\setAn{1}, \ldots, \setAn{m}, \setBn{1}, \ldots, \setBn{n}}
\end{equation}
for any $n, m \setin \natnumbers$.

The elements of~$\cObj{\setAn{1}, \ldots, \setAn{m}} \cprod \cObj{\setBn{1}, \ldots, \setBn{n}}$ are all possible concatenations of elements of~$\cObj{\setAn{1}, \ldots, \setAn{m}}$ with elements of~$\cObj{\setBn{1}, \ldots, \setBn{n}}$.

Equation \cref{eq:strict-monoidal-product-sets} holds in particular also in the cases when~$m = 0$ or~$n =0$:
\begin{equation}
    \makecprod{
        \emptycobj,
        \cObj{\setBn{1}, \ldots, \setBn{n}}
    } =
    \cObj{
        \setBn{1}, \ldots, \setBn{n}
    }
\end{equation}
and
\begin{equation}
    \makecprod{
        \cObj{\setAn{1}, \ldots, \setAn{m}},
        \emptycobj
    } =
    \cObj{
        \setAn{1}, \ldots, \setAn{m}
    }.
\end{equation}

% Also, writing $\Emptyset = \cObj{\Emptyset}$ we find that
% \begin{equation}
%     \makecprod{
%         \cObj{\setAn{1}, \ldots, \setAn{m}},
%         \Emptyset
%     } = \Emptyset
% \end{equation}
% and
% \begin{equation}
%     \makecprod{\Emptyset, \cObj{\setBn{1}, \ldots, \setBn{n}}} = \Emptyset.
% \end{equation}

The multiplication of objects in~\SetL is \SY{associative} (which was our goal with making this construction).
Indeed, both
\begin{equation}
    \makecprod{
        ( \makecprod{\cObj{\setAn{1}, \ldots, \setAn{l}}, \cObj{\setBn{1}, \ldots, \setBn{m}}})
        ,
        \cObj{\setCn{1}, \ldots, \setCn{n}}
    }
\end{equation}
and
\begin{equation}
    \makecprod{
        \cObj{\setAn{1}, \ldots, \setAn{l}},
        (
        \makecprod{
            \cObj{\setBn{1}, \ldots, \setBn{m}},
            \cObj{\setCn{1}, \ldots, \setCn{n}}
        }
        )
    }
\end{equation}
are equal to
\begin{equation}
    \cObj{\setAn{1}, \ldots, \setAn{l}, \setBn{1}, \ldots, \setBn{m}, \setCn{1}, \ldots, \setCn{n}}.
\end{equation}

\begin{remark}
    Note that, for any sets $\setA, \setB$, we have
    \begin{equation}
        \cObj{\setA, \setB} = \setA \cartprod \setB.
    \end{equation}
    In other words, the multiplication operation in $\SetL$ ``coincides'' with the cartesian product operation in the special case when we are multiplying two sets together.

    Because of this, the operation of cartesian product is well-defined for $\SetL$.
    Namely, given tuples sets $\cObj{\setAn{1}, \ldots, \setAn{l}}$ and $\cObj{\setBn{1}, \ldots, \setBn{m}}$, their cartesian product is again a tuple set:
    \begin{equation}
        \cObj{\setAn{1}, \ldots, \setAn{l}} \cartprod \cObj{\setBn{1}, \ldots, \setBn{m}} = \cObj{\cObj{\setAn{1}, \ldots, \setAn{l}}, \cObj{\setBn{1}, \ldots, \setBn{m}}},
    \end{equation}
    (which is a ``nested'' tuple set).

\end{remark}

\showslides{
    \begin{forslides}
        \begin{equation}
            % \label{eq:moore-assoc-a-bis}
            \mora \mthen (\morb \mthen \morc)
        \end{equation}
        %
        \begin{equation}
            % \label{eq:moore-assoc-b-bis}
            (\mora \mthen \morb) \mthen \morc
        \end{equation}
        %
        \begin{equation}
            % \label{eq:lset}
            \SetL
        \end{equation}
        %
        \begin{equation}
            % \label{eq:definition-setl}
            \setAn{1}, \setAn{2}, \ldots, \setAn{n}
        \end{equation}
        %
        \begin{equation}
            % \label{eq:definition-setlbis}
            n\setin \natnumbers
        \end{equation}
        \begin{equation}
            % \label{eq:mult-setl}
            \cObj{\setAn{1}, \setAn{2}, \ldots, \setAn{m}}
        \end{equation}
        \begin{equation}
            % \label{eq:mult-setlb}
            \cObj{\setBn{1}, \setBn{2}, \ldots, \setBn{n}}
        \end{equation}
    \end{forslides}
}


\begin{lemma}
    \label{lem:SetL-is-associative-stacking}
    \SetL is  \SY{associative stacking} using the structure that arises from tuple concatenation.
    % by setting $\mtimescatob_{\SetL} \definedas \cprod$
    % and $\mtimescatmor_{\SetL} \definedas \cprod$
\end{lemma}
\begin{proof}
    % In this example we show that~\SetL is an  \SY{associative stacking} semicategory.
    % We first show that it is a stacking category.
    % To do so, we start by showing that the stacking operations fulfill \cref{def:simple-stacking-semi-cat}.
    For the stacking operation on objects $\mtimescatob$ we use the operation $\cprod$ defined in \cref{sec:cartcatset} which was referred to as the ``multiplication in \SetL'':
    \begin{align}
        \cObj{\setAn{1}, \dots, \setAn{m}} \mtimescatob \cObj{\setBn{1}, \dots, \setBn{n}} & \definedas \cObj{\setAn{1}, \dots, \setAn{m}} \cprod \cObj{\setBn{1}, \dots, \setBn{n}} \\
                                                                                           & = \cObj{\setAn{1}, \dots, \setAn{m}, \setBn{1}, \dots, \setBn{n}}.
    \end{align}
    It was shown there that this operation is associative.

    As for $\mtimescatmor$, we define it as follows:
    \begin{equation}
        \label{eq:mtimescatmor-SetL}
        \prfperiod{
            \mora\colon\setA\mtoin\SetL\setB
        }{
            \quad
        }{
            \morb\colon\setC\mtoin\SetL\setD
        }{
            \defmap{
                (\mora\mtimescatmor\morb )
            }{
                \setA \cprod \setC
            }{
                \mtoin\SetL
            }{
                \setB \cprod \setD
            }{
                \maketupconcat{\ela, \elc}
            }{
                \maketupconcat{\mora(\ela), \morb(\elc)}
            }
        }
    \end{equation}
    The two operations $\mtimescatob, \mtimescatmor$ so defined satisfy the compatibility conditions
    required by \cref{def:simple-stacking-semi-cat}.

    To show associativity, consider three morphisms
    \begin{equation}\label{eq:SetL-morphisms-to-associatively-stack}
        % \begin{aligned}
        \mora\colon\setA\mtoin\SetL\setB,\qquad
        \morb\colon\setC\mtoin\SetL\setD,\qquad
        \morc\colon\setE\mtoin\SetL\setF.
        % \end{aligned}
    \end{equation}
    We compute $\mora\mtimescatmor(\morb\mtimescatmor\morc)$ and $(\mora\mtimescatmor\morb)\mtimescatmor\morc$ following the recipe \cref{eq:mtimescatmor-SetL} to obtain
    \begin{equation}\label{eq:SetL-morphisms-associative-stack-check-1}
        \defmapcomma{
            \parslight{\mora\mtimescatmor\morb}\mtimescatmor\morc
        }{
            \parslight{\setA \cprod \setC} \cprod \setE
        }{
            \mtoin\SetL
        }{
            \parslight{\setB \cprod \setD} \cprod \setF
        }{
            \maketupconcat{\parslight{\maketupconcat{\ela, \elc}}, \ele}
        }{
            \maketupconcat{\parslight{\maketupconcat{\mora(\ela), \morb(\elc)}}, \morc(\ele) }
        }
    \end{equation}
    \begin{equation}\label{eq:SetL-morphisms-associative-stack-check-2}
        \defmapperiod{
            \mora\mtimescatmor\parslight{\morb\mtimescatmor\morc}
        }{
            \setA \cprod \parslight{\setC \cprod \setE}
        }{
            \mtoin\SetL
        }{
            \setB \cprod \parslight{\setD \cprod \setF}
        }{
            \maketupconcat{\ela, \parslight{\maketupconcat{\elc, \ele}}}
        }{
            \maketupconcat{\mora(\ela), \parslight{\maketupconcat{\morb(\elc), \morc(\ele)}}}
        }
    \end{equation}
    Notice that the operations $\cprod$ and $\tupconcat$ are associative; therefore, we can remove all the light parentheses that appear in the formulas.
    This implies that both functions are equal to
    \begin{equation}\label{eq:SetL-morphisms-associative-stack-check-3}
        \defmapperiod{
            \mora\mtimescatmor\morb\mtimescatmor\morc
        }{
            \setA \cprod \setC \cprod \setE
        }{
            \mtoin\SetL
        }{
            \setB \cprod \setD \cprod \setF
        }{
            \maketupconcat{\ela, \elc, \ele}
        }{
            \maketupconcat{\mora(\ela), \morb(\elc), \morc(\ele)}
        }
    \end{equation}
    % \begin{equation}
    %     \defmapperiod{
    %         \mtimescatob
    %     }{
    %         \Obof{\SetL} \cartprod \Obof{\SetL}
    %     }{
    %         \mto
    %     }{
    %         \Obof{\SetL}
    %     }{
    %         \tup{\cObj{\setAn{1},\ldots, \setAn{n}}, \cObj{\setBn{1},\ldots,\setBn{m}}}
    %     }{
    %         \cObj{\setAn{1},\ldots, \setAn{n}, \setBn{1},\ldots,\setBn{m}}
    %     }
    % \end{equation}
    % The stacking operation on morphisms is defined as
    % \begin{equation}
    %     \defmapcomma{
    %         \mtimescatmor
    %     }{
    %         \Mor_{\SetL} \cartprod \Mor_{\SetL}
    %     }{
    %         \mto
    %     }{
    %         \Mor_{\SetL}
    %     }{
    %         \tup{\mora,\morb}
    %     }{
    %         \mora \cprod \morb
    %     },
    % \end{equation}
    % where~$(\mora \cprod \morb)(\tup{\setAeln{1},\ldots, \setAeln{m}} \tupconcat \tup{\setBel_1,\ldots, \setBel_n})\definedas \mora(\tup{\setAeln{1},\ldots, \setAeln{m}})\tupconcat \morb(\tup{\setBel_1,\ldots, \setBel_n})$.
    % This is indeed a map between sets.
    % The two operations are compatible, because
    % \begin{widepar}
    %     \begin{equation}
    %         \prfperiod{
    %             \mora\colon \cObj{\setAn{1},\ldots,\setAn{n}}\mto \cObj{\setBn{1},\ldots,\setBn{m}}
    %         }{
    %             \morb\colon \cObj{\setCn{1},\ldots,\setCn{o}}\mto \cObj{\setDn{1},\ldots,\setDn{p}}
    %         }{
    %             \mora \mtimescatmor \morb \colon \cObj{\setAn{1}, \ldots, \setAn{n}, \setCn{1}, \ldots, \setCn{o}} \mto \cObj{\setBn{1}, \ldots, \setBn{m}, \setDn{1}, \ldots, \setDn{p}}
    %         }
    %     \end{equation}
    % \end{widepar}

    % We now show associativity.
    % We have already shown in \cref{sec:cartcatset} that~$\cprod$ satisfies the properties of a \SY{semigroup} operation.
    % For the operation on morphisms, consider
    % \begin{equation}
    %     \begin{aligned}
    %         \mora\colon \cObj{\setAn{1},\ldots,\setAn{m}} & \mto \cObj{\setBn{1}, \ldots, \setBn{n}}, \\
    %         \morb\colon \cObj{\setCn{1},\ldots,\setCn{o}} & \mto \cObj{\setDn{1},\ldots,\setDn{p}}, \\
    %         \morc\colon \cObj{\setEn{1},\ldots,\setEn{q}} & \mto \cObj{\setFn{1},\ldots,\setFn{r}},
    %     \end{aligned}
    % \end{equation}
    % For brevity, we write~$\setAeln{}=\tup{\setAeln{1},\ldots, \setAeln{m}}$,~$\setCeln{}=\tup{\setCeln{1},\ldots, \setCeln{o}}$, and~$\setEeln{}=\tup{\setEeln{1},\ldots,\setEeln{q}}$.
    % We have
    % \begin{align}
    %     ((\mora \cprod \morb)\cprod \morc)(\tup{\setAeln{}, \setCeln{},\setEeln{}})
    %      & =(\mora(\setAeln{})\tupconcat \morb(\setCeln{}))\tupconcat \morc(\setEeln{}) \\
    %      & =\mora(\setAeln{})\tupconcat \morb(\setCeln{})\tupconcat \morc(\setEeln{}) \\
    %      & =\mora(\setAeln{})\tupconcat (\morb(\setCeln{})\tupconcat \morc(\setEeln{})) \\
    %      & =(\mora \cprod (\morb\cprod \morc))(\tup{\setAeln{}, \setCeln{},\setEeln{}})
    % \end{align}
    % for all~$\tup{\setAeln{}, \setCeln{},\setEeln{}}\setin \cObj{\setAn{1},\ldots,\setAn{m},\setCn{1},\ldots,\setCn{o},\setEn{1},\ldots,\setEn{q}}$.

\end{proof}

% \begin{comment}
% \todotext{The example below is perhaps deprecated if we no longer use Set* as a main character in our story}

% \begin{example}
%     \SetL is an  \SY{associative stacking} semicategory.
%     We first show that it is a stacking category.
%     The stacking operation on objects is defined as list concatenation:
%     \begin{equation}
%         \defmapperiod{\mtimescatob}{\Ob_{\SetL} \cartprod \Ob_{\SetL}}{\mto}{\Ob_{\SetL}}{\tup{\Tupcatt{\setAn{1}}{\ldots}{ \setAn{n}}, \Tupcatt{\setBn{1}}{\ldots}{\setBn{m}}}}{\Tupcatt{\setAn{1}}{\ldots}{\setAn{n}} \listconcat \Tupcatt{\setBn{1}}{\ldots}{\setBn{m}}}
%     \end{equation}
%     The stacking operation on morphisms is defined as
%     \begin{equation}
%         \defmapperiod{\mtimescatmor}{\Mor_{\SetL} \cartprod \Mor_{\SetL}}{\mto}{\Mor_{\SetL}}{\tup{\mora,\morb}}{\more},
%     \end{equation}
%     which, given~$\mora\colon \Tupcatt{\setAn{1}}{\ldots}{\setAn{n}}\mto \Tupcatt{\setBn{1}}{\ldots}{\setBn{m}}$ (given by a map~$\mapc$) and~$\Tupcatt{\setCn{1}}{\ldots}{\setCn{o}}\mto \Tupcatt{\setDn{1}}{\ldots}{\setDn{p}}$ (given by a map~$\mapd$) is given by a function
%     \begin{equation}
%         \mape \definedas \cohm_{\setA,\setC}\then (\mapc\funcprod \mapd) \then \cohm_{\setB,\setD},
%     \end{equation}
%     which has signature
%     \begin{widepar}
%         \begin{equation}
%             \setAn{1}\cartprod (\setA_2\cartprod (\ldots \cartprod (\setAn{n} \cartprod (\setCn{1} \cartprod(\setC_2 \cartprod (\ldots \cartprod \setCn{o}))))))
%             \to
%             \setBn{1}\cartprod (\setB_2\cartprod (\ldots \cartprod (\setBn{m} \cartprod (\setDn{1} \cartprod (\setD_2 \cartprod(\ldots \cartprod \setDn{p})))))).
%         \end{equation}
%     \end{widepar}
%     The two operations are compatible, because
%     \begin{widepar}
%         \begin{equation}
%             \prfperiod{\mora\colon \Tupcatt{\setAn{1}}{\ldots}{\setAn{n}}\mto \Tupcatt{\setBn{1}}{\ldots}{\setBn{m}}}
%             {\morb\colon \Tupcatt{\setCn{1}}{\ldots}{\setCn{o}}\mto \Tupcatt{\setDn{1}}{\ldots}{\setDn{p}}}
%             {\more\colon \Tupca{\setAn{1}\setconcat \ldots \setconcat \setAn{n} \setconcat \setCn{1} \setconcat \ldots \setconcat \setCn{o}} \mto \Tupca{\setBn{1}\setconcat \ldots \setconcat \setBn{m} \setconcat \setDn{1} \setconcat \ldots \setconcat \setDn{p}}}
%         \end{equation}
%     \end{widepar}
%     Note that using the coherence morphisms is needed for the compatibility conditions to hold.

%     We now show associativity.
%     We have already shown in \cref{sec:semigroups} that the concatenation of lists satisfies the properties of a \SY{semigroup} operation.
%     Showing associativity of the operation on the morphisms is more cumbersome.
%     Consider morphisms
%     \begin{equation}
%         \begin{aligned}
%             \mora\colon \Tupcatt{\setAn{1}}{\ldots}{\setAn{m}} & \mto \Tupcatt{\setBn{1}}{\ldots}{\setBn{n}}, \\
%             \morb\colon \Tupcatt{\setCn{1}}{\ldots}{\setCn{o}} & \mto \Tupcatt{\setDn{1}}{\ldots}{\setDn{p}}, \\
%             \morc\colon \Tupcatt{\setEn{1}}{\ldots}{ \\setEn{q}} & \mto \Tupcatt{\\setFn{1}}{\ldots}{\\setFn{r}},
%         \end{aligned}
%     \end{equation}
%     to which maps~$\mapd,\mape,\mapf$ are associated.

%     We have:

%     \begin{center}
%         \begin{tikzpicture}
%             \node at (0,0){
%                 \begin{tikzcd}[every arrow/.append style={-Triangle, draw=morphisms}, column sep=large]
%                     (\mora \mtimescatmor \morb)
%                     \mtimescatmor \morc                                                                                                                                                                      & \mora \mtimescatmor (\morb \mtimescatmor \morc) \\
%                     \tupset{((\setLA \listconcat \setLC)\listconcat \setLE)}\arrow[d,"\cohm_{\setLA\listconcat \setLC, \setLE}^{-1}"]\arrow[r, equal]                                                        & \tupset{(\setLA \listconcat (\setLC\listconcat \setLE))}\arrow[d] \\
%                     \tupset{(\setLA\listconcat \setLC)}\cartprod \tupset{(\setLE)}\arrow[d, "\cohm_{\tupset{(\setLA)}, \tupset{(\setLC)}}^{-1} \funcprod \catid_{\tupset{(\setLC)}}"]                         & \tupset{(\setLA)}\cartprod \tupset{(\setLC \listconcat \setLE)}\arrow[d] \\
%                     (\tupset{(\setLA)}\cartprod \tupset{(\setLC)})\cartprod \tupset{(\setLE)}\arrow[r, "\alpha"]\arrow[d, "(\mapd \funcprod \mape)\funcprod \mapf"]                                          & \tupset{(\setLA)}\cartprod (\tupset{(\setLC)}\cartprod \tupset{(\setLE)})\arrow[d, "\mapd \funcprod (\mape\funcprod \mapf)"] \\
%                     (\tupset{(\setLB)}\cartprod \tupset{(\setLD)})\cartprod \tupset{(\setLF)}\arrow[r, "\alpha"]\arrow[d, "\cohm_{\tupset{(\setLB)}, \tupset{(\setLD)}}\funcprod \catid_{\tupset{(\setLF)}}"] & \tupset{(\setLB)}\cartprod (\tupset{(\setLD)}\cartprod \tupset{(\setLF)})\arrow[d] \\
%                     \tupset{(\setLB\listconcat \setLD)}\cartprod \tupset{(\setLF)}\arrow[d,"\cohm_{\setLB\listconcat \setLD, \setLF}"]                                                                       & \tupset{(\setLB)} \cartprod \tupset{(\setLD \listconcat\setLF)}\arrow[d] \\
%                     \tupset{((\setLB \listconcat \setLD)\listconcat \setLF)}\arrow[r, equal]                                                                                                                 & \tupset{(\setLB \listconcat (\setLD\listconcat \setLF))}
%                 \end{tikzcd}
%             };
%         \end{tikzpicture}
%     \end{center}
%     \todo{finish diagram above}
% \end{example}
% \end{comment}

%
%\begin{definition}
%    An  \SY{associative stacking} semicategory is called \emph{commutative} if its two stacking operations make the collections of objects and morphisms, respectively, into commutative semigroups.
%\end{definition}


\subsection{Moore machines:~\SetL version}
We can now slightly rework the definition of \SY{Moore machines} using morphisms in \SetL rather than \Set.
\begin{definition}[Moore machine, \SetL version]
    \label{def:moore-machine}
    A \maindef{Moore machine} is a tuple
    \begin{equation}
        \label{eq:moore-tuple-improved-again}
        \tup{\prinL,\prstL,\proutL,\prdyn,\prreadout,\prstart},
    \end{equation}
    where~$\prinL$,~$\prstL$, and~$\proutL$ are objects of~$\SetL$,
    \begin{equation}
        \prdyn \colon \prinL \cprod \prstL \mto \prstL,
    \end{equation}
    and
    \begin{equation}
        \prreadout \colon \prstL \mto \proutL,
    \end{equation}
    are morphisms in $\SetL$ (so they are functions),
    and~$\prstart \setin \prstL$.
\end{definition}

\paragraph{Composition}
\label{sec:composition-of-Moore-machines}

Consider two Moore machines
\begin{equation}
    \label{eq:moore-mora-3rd}
    \mora = \tupp{\prinL_{\mora},\prstL_{\mora},\proutL_{\mora},\prdyn_{\mora},\prreadout_{\mora},\prstart_{\mora}}
\end{equation}
\begin{equation}
    \label{eq:moore-morb-3rd}
    \morb = \tupp{\prinL_{\morb},\prstL_{\morb},\proutL_{\morb},\prdyn_{\morb},\prreadout_{\morb},\prstart_{\morb}},
\end{equation}
that are compatible for composition, in the sense that $\proutL_{\mora} = \prinL_{\morb}$.
Their composition is given by
\begin{equation}\label{eq:moore-setL-composite}
    \morab = \tupp{\prinL_{\morab},\prstL_{\morab},\proutL_{\morab},\prdyn_{\morab},\prreadout_{\morab},\prstart_{\morab}},
\end{equation}
where
%
\begin{equation}
    \label{eq:moore-comp-3rd-1}
    \begin{aligned}
        \prinL_{\morab}   & \definedas \prinL_{\mora}, \\
        \prstL_{\morab}   & \definedas \prstL_{\mora} \cprod \prstL_{\morb}, \\
        \proutL_{\morab}  & \definedas \proutL_{\morb}, \\
        \prstart_{\morab} & \definedas \prstart_{\mora} \tupconcat \prstart_{\morb},
    \end{aligned}
\end{equation}
and~$\prdyn_{\morab}$ and~$\prreadout_{\morab}$ are defined as
%
\begin{equation}
    \label{eq:moore-comp-3rd-2}
    \defmapcommaset{
        \prdyn_{\morab}
    }{
        \prinL_{\mora} \cprod \prstL_{\mora} \cprod \prstL_{\morb}
    }{
        \prstL_{\mora} \cprod \prstL_{\morb}
    }{
        \prinel \tupconcat \prstel_{\mora} \tupconcat \prstel_{\morb}
    }{
        \prdyn_{\mora} (\prinel \tupconcat \prstel_{\mora}) \tupconcat \prdyn_{\morb}(\prreadout_{\mora}(\prstel_{\mora}) \tupconcat \prstel_{\morb})
    }
\end{equation}
and
\begin{equation}
    \label{eq:moore-comp-3rd-3}
    \defmapperiodset{
        \prreadout_{\morab}
    }{
        \prstL_{\mora} \cprod \prstL_{\morb}
    }{
        \proutL_{\morb}
    }{
        \prstel_{\mora} \tupconcat \prstel_{\morb}
    }{
        \prreadout_{\morb}(\prstel_{\morb})
    }
\end{equation}

\paragraph{Composition is associative}

Let three composable \SY{Moore machines}~$\mora$,~$\morb$, and~$\morc$ be given.
We check that each of the six entries in the definition \cref{def:moore-machine} coincide for~$(\morab) \mthen \morc$ and~$\mora \mthen (\morb \mthen \morc)$.

Clearly,
\begin{equation}
    \label{eq:moore_assoc_input}
    \prinL_{(\morab)\mthen \morc} = \prinL_{\mora} =\prinL_{\mora \mthen (\morb \mthen \morc)}
\end{equation}
and
\begin{equation}
    \label{eq:moore_assoc_output}
    \proutL_{(\morab)\mthen \morc} =\proutL_\morc = \proutL_{\mora \mthen (\morb \mthen \morc)}.
\end{equation}
Furthermore,
\begin{equation}
    \label{eq:moore_assoc_state}
    \prstL_{(\morab)\mthen \morc} = (\prstL_{\mora} \cprod \prstL_{\morb}) \cprod \prstL_{\morc} = \prstL_{\mora} \cprod (\prstL_{\morb} \cprod \prstL_{\morc}) =\prstL_{\mora \mthen (\morb \mthen \morc)}
\end{equation}
since concatenation of lists is \SY{associative}.
Similarly,
\begin{equation}
    \label{eq:moore_assoc_start}
    \begin{aligned}
        \prstart_{(\morab)\mthen \morc} & = \prstart_{\morab} \tupconcat \prstart_{\morc} \\
                                        & = (\prstart_{\mora } \tupconcat \prstart_{\morb}) \tupconcat \prstart_{\morc} \\
                                        & = \prstart_{\mora } \tupconcat (\prstart_{\morb} \tupconcat \prstart_{\morc}) \\
                                        & = \prstart_{\mora } \tupconcat \prstart_{ \morb \mthen\morc} \\
                                        & = \prstart_{\mora \mthen (\morb \mthen \morc)}.
    \end{aligned}
\end{equation}
Next we show that~$\prdyn_{(\mora \mthen \morb)\mthen \morc} = \prdyn_{\mora \mthen (\morb \mthen \morc)}$.

On the one hand,
\begin{widepar}
    \begin{equation}
        \label{eq:assoc_moore_3rd_1}
        \begin{aligned}
            \prdyn_{(\morab)\mthen \morc}\colon \prinL_{\mora} \cprod \prstL_{\mora} \cprod \prstL_{\morb} \cprod \prstL_{\morc} & \to \prstL_{\mora} \cprod \prstL_{\morb} \cprod \prstL_{\morc} \\
            \prinel \tupconcat \prstel_{\mora} \tupconcat \prstel_{\morb} \tupconcat \prstel_\morc                               & \mapsto \prdyn_{\morab}(\prinel \tupconcat \prstel_\mora \tupconcat \prstel_\morb) \tupconcat \prdyn_{\morc}(\prreadout_{\morab}(\prstel_\mora \tupconcat \prstel_\morb) \tupconcat \prstel_\morc) \\
                                                                                                                                 & = \prdyn_\mora(\prinel \tupconcat \prstel_\mora) \tupconcat \prdyn_\morb(\prreadout_\mora(\prstel_\mora) \tupconcat \prstel_\morb) \tupconcat \prdyn_\morc(\prreadout_\morb(\prstel_\morb) \tupconcat \prstel_\morc),
        \end{aligned}
    \end{equation}
\end{widepar}
while on the other hand
\begin{widepar}
    \begin{equation}
        \label{eq:assoc_moore_3rd_2}
        \begin{aligned}
            \prdyn_{\mora \mthen (\morb \mthen \morc)} \colon \prinL_{\mora} \cprod \prstL_{\mora} \cprod \prstL_{\morb} \cprod \prstL_{\morc} & \to \prstL_{\mora} \cprod \prstL_{\morb} \cprod \prstL_{\morc} \\
            \prinel \tupconcat \prstel_{\mora} \tupconcat \prstel_{\morb} \tupconcat \prstel_\morc                                             & \mapsto \prdyn_{\mora }(\prinel \tupconcat \prstel_\mora) \tupconcat \prdyn_{\morb \mthen \morc}(\prreadout_{\mora }(\prstel_\mora) \tupconcat \prstel_\morb \tupconcat \prstel_\morc) \\
                                                                                                                                               & = \prdyn_\mora(\prinel \tupconcat \prstel_\mora) \tupconcat \prdyn_\morb(\prreadout_\mora(\prstel_\mora) \tupconcat \prstel_\morb) \tupconcat \prdyn_\morc(\prreadout_\morb(\prstel_\morb) \tupconcat \prstel_\morc).
        \end{aligned}
    \end{equation}
\end{widepar}
So, these two functions are indeed the same.

Finally, we verify that~$\prreadout_{(\morab) \mthen \morc} = \prreadout_{\mora\mthen (\morb \mthen \morc)}$:
\begin{equation}
    \label{eq:assoc_moore_3rd_3}
    \defmapset{
        \prreadout_{(\morab) \mthen \morc}
    }{
        \prstL_{\mora} \cprod \prstL_{\morb} \cprod \prstL_{\morc}
    }{
        \proutL_{\morc}
    }{
        \prstel_{\mora} \tupconcat \prstel_{\morb} \tupconcat \prstel_{\morc}
    }{
        \prreadout_{\morc}(\prstel_{\morc})
    }
\end{equation}
while
\begin{equation}
    \label{eq:assoc_moore_3rd_4}
    \defmapperiodset{
        \prreadout_{\mora \mthen (\morb \mthen \morc)}
    }{
        \prstL_{\mora} \cprod \prstL_{\morb} \cprod \prstL_{\morc}
    }{
        \proutL_{\morc}
    }{
        \prstel_{\mora} \tupconcat \prstel_{\morb} \tupconcat \prstel_{\morc}
    }{
        \prreadout_{\morb \mthen \morc}(\prstel_{\morb} \tupconcat \prstel_{\morc}) = \prreadout_{\morc}(\prstel_{\morc})
    }
\end{equation}

\paragraph{The \SY{semicategory} of Moore machines}

Now that we have shown that composition of \SY{Moore machines} is \SY{associative} (with our new definition), we can organize \SY{Moore machines} as a \SY{semicategory}.

\begin{definition}[\Moore]
    \label{def:Moore}
    The \maindef{semicategory of Moore machines} \Moore is given by:
    \begin{enumerate}
        \item \emph{Objects:} objects of~\SetL.
        \item \emph{Morphisms:}
              A morphism from~$\prinL$ to $\proutL$ is a Moore machine
              \begin{equation}
                  \tup{\prinL,\prstL,\proutL,\prdyn,\prreadout,\prstart},
              \end{equation}
              where
              \begin{itemize}
                  \item $\prinL,\prstL,\proutL$ are objects of \SetL;
                  \item $\prdyn \colon \prinL \cprod \prstL \mto \prstL$ and $\prreadout \colon \prstL \mto \proutL$ are morphisms in $\SetL$;
                  \item $\prstart \setin \prstL$.
              \end{itemize}
        \item \emph{Composition:}
              as defined above in this section.
    \end{enumerate}
\end{definition}

\todotext{\alphubel: Add a remark showing that \SY{Moore machines} do not have identities, and are hence do not form a category.}


\subsection{The category~\RelL}
We define a category~\RelL where the objects are sets of tuples (as in~\SetL).

\begin{definition}[The category~\RelL]
    \label{def:RelL}
    \SYNDEF{category of tuple-sets and relations}
    The category~\RelL is the \SY{subcategory} of~\Rel where the objects are tuples sets (objects of \SetL).
    % \begin{enumerate}
    %     \item \emph{Objects:} sets of tuples (as in~\SetL)
    %     \item \emph{Morphisms:}
    %           relations between sets of tuples.
    %     \item \emph{Composition:}
    %           Composition is the usual composition of relations.
    %     \item \emph{Identities:}
    %           The identity morphism on an object~$\cObj{\setAn{1},\ldots, \setAn{m}}$ is given by the identity relation~$\catid_{\cObj{\setAn{1},\ldots, \setAn{m}}}$.
    % \end{enumerate}
\end{definition}

\begin{lemma}\label{lem:RelL-associative-stacking}
    \RelL is  \SY{associative stacking} with the structure induced by tuple concatenation.
\end{lemma}
\begin{proof}
    The multiplication $\mtimescatob$ is $\cprod$, the same as the one defined for $\SetL$.

    The multiplication $\mtimescatmor$ is defined as follows:
    \begin{equation}
        \prf{
            \relA\setsubseteq \setA \cartprod \setB
        }{
            \quad
        }{
            \relB\setsubseteq \setC\cartprod \setD
        }{
            \pars{\relA\mtimescatmor\relB} \setsubseteq
            \pars{\makecprod{\setA,\setC}} \cartprod
            \pars{\makecprod{\setB,\setD}}
        }
    \end{equation}
    where
    \begin{equation}
        \pars{
            \relA
            \mtimescatmor
            \relB
        }
        =
        \makesett{
            \tup{\setAeln{}\tupconcat \setCeln{},\setBeln{}\tupconcat \setDeln{}}
            \mid
            % \setAeln{} \setin \setA,
            % \setBeln{} \setin \setB,
            % \setCeln{} \setin \setC,
            % \setDeln{} \setin \setD:
            (\inrel{\setAeln{}}\relA{\setBeln{}})
            \booland
            (\inrel{\setCeln{}}\relB{\setDeln{}})
        }.
    \end{equation}
    The rest of the proof is left as an exercise.
\end{proof}

% What is a relation in this context?
% Given two sets of tuples~$\cObj{\setAn{1},\ldots, \setAn{m}}$ and~$\cObj{\setBn{1},\ldots, \setBn{n}}$, a relation between them is given by
% \begin{equation}
%     \relA \setsubseteq \cObj{\setAn{1},\ldots, \setAn{m}} \cartprod \cObj{\setBn{1},\ldots, \setBn{n}}.
% % \end{equation}
% We can define a multiplication operation for relations.
% Consider
% \begin{equation}
%     \begin{aligned}
%         \relA\colon \cObj{\setAn{1},\ldots, \setAn{m}} & \mtoin{\RelL}\cObj{\setBn{1},\ldots, \setBn{n}}, \\
%         \relB\colon \cObj{\setCn{1},\ldots, \setCn{o}} & \mtoin{\RelL}\cObj{\setDn{1},\ldots, \setDn{p}}.
%     \end{aligned}
% \end{equation}
% We have
% \begin{equation}
%     \relA\mtimescatmor \relB \colon \cObj{\setAn{1},\ldots, \setAn{m}}\cprod \cObj{\setCn{1},\ldots, \setCn{o}}\mtoin{\RelL}
%     \cObj{\setBn{1},\ldots, \setBn{n}}\cprod \cObj{\setDn{1},\ldots, \setDn{p}}.
% \end{equation}
% Specifically, the multiplication is given by:
% \begin{equation}
%     \label{eq:multiofrels}
%     \defmapperiod{
%         \mtimescatmor
%     }{
%         \Mor_{\RelL}\cartprod \Mor_{\RelL}
%     }{
%         \mto
%     }{
%         \Mor_{\RelL}
%     }{
%         \tup{\relA,\relB}
%     }{
%         \makeset{
%             \tup{\setAeln{}\tupconcat \setCeln{},\setBeln{}\tupconcat \setDeln{}}
%             \mid
%             (\inrel{\setAeln{}}\relA{\setBeln{}}) \booland (\inrel{\setCeln{}}\relB{\setDeln{}})
%         }
%     }
% \end{equation}

\begin{exercise}
    \label{ex:relassocstack}
    Consider the stacking operations on objects and on morphisms introduced in this section.
    Prove that~\RelL is \SY{associative stacking}.
\end{exercise}

\begin{solution}
    To prove the statement we first check that the stacking operations satisfy \cref{def:simple-stacking-semi-cat}, we then show that they are compatible, and finally show associativity.
    The stacking operation on objects was already checked for~\SetL.
    The stacking operation on morphisms clearly returns a valid relation.
    Furthermore, compatibility is satisfied:
    \begin{equation}
        \prfperiod{
            \relA\colon \cObj{\setAn{1},\ldots, \setAn{m}} \mtoin{\RelL}\cObj{\setBn{1},\ldots, \setBn{n}}
        }
        {
            \relB\colon \cObj{\setCn{1},\ldots, \setCn{o}} \mtoin{\RelL}\cObj{\setDn{1},\ldots, \setDn{p}}
        }
        {
            \relA\mtimescatmor \relB \colon \cObj{\setAn{1},\ldots, \setAn{m}}\cprod \cObj{\setCn{1},\ldots, \setCn{o}}\mtoin{\RelL}
            \cObj{\setBn{1},\ldots, \setBn{n}}\cprod \cObj{\setDn{1},\ldots, \setDn{p}}
        }
    \end{equation}
    Finally, associativity for the operation on objects was already shown for~\SetL.
    % The stacking operation on morphisms follows directly from \cref{eq:multiofrels}.
\end{solution}

\begin{ctdefinition}[Trace for $\RelL$ with concatenation as monoidal product]
    \label{def:tuple-REL-trace}
    Consider the symmetric strict monoidal category $\RelL$, where the monoidal product is defined via the concatenation operation $\cprod$ for lists of sets.
    Given a relation $\relA \colon \cObj{\setA, \setC}  \to \cObj{\setB, \setC}$, its \emph{trace} is defined as
    %
    \begin{equation}
        \label{eq:tupel-REL-tracedef}
        \begin{aligned}
            \Tr_{\cObj{\setA}, \cObj{\setB}}^{\cObj{\setC}}(\relA) = \{ \tup{\ela, \elb} \in \cObj{\setA} \times \cObj{\setB} \mid \exists \ \elc \in \cObj{\setC} : \tup{\ela \tupconcat \elc, \elb \tupconcat \elc} \in \relA \}
        \end{aligned}
    \end{equation}
\end{ctdefinition}


\subsection{The category~\PosL}

We define a category analogous to~\SetL, but its objects are ``tuple posets''.

Given \SY{posets}~$\posA_1,\ldots, \posA_n$ we define the poset
\begin{equation}
    \label{eq:posetoftuples}
    \cObj{\posA_1,\ldots, \posA_n}\definedas \tup{\cObj{\posAsetn{1},\ldots, \posAsetn{n}}, \posleqof{\cObj{\posA_1,\ldots, \posA_n}}},
\end{equation}
where~$\cObj{\posAsetn{1},\ldots, \posAsetn{n}}$ is a set of tuples, and we use the product order:
\begin{equation}
    \prfdoubleperiod{
        \tup{\posela_1,\ldots, \posela_n}
        \posleqof{\cObj{\posA_1,\ldots, \posA_n}}
        \tup{\poselb_1,\ldots, \poselb_n}
    }{
        \posela_i \posleqof{\posA_i} \poselb_i \text{ for all }~i\setin \makeset{1,\ldots, n}
    }
\end{equation}
\begin{definition}[The category~\PosL]
    \SYNDEF{category of tuple-posets and monotone maps}
    \label{def:PosL}
    The category~\PosL is the \SY{subcategory} of \Pos where the objects are tuple posets, \SY{posets} of the form
    \begin{equation}
        \cObj{\posA_1,\ldots, \posA_n} = \tup{\cObj{\posAsetn{1},\ldots, \posAsetn{n}}, \posleqof{\cObj{\posA_1,\ldots, \posA_n}}}.
    \end{equation}
    % \begin{enumerate}
    %     \item \emph{Objects:} \SY{posets} of the form presented in \cref{eq:posetoftuples}.
    %     \item \emph{Morphisms:}
    %           Monotone functions between \SY{posets} of the form presented in \cref{eq:posetoftuples}.
    %     \item \emph{Composition:}
    %           Composition is the usual composition of functions.
    %     \item \emph{Identities:}
    %           Identity functions.
    % \end{enumerate}
\end{definition}

\begin{lemma}\label{lem:PosL-associative-stacking}
    \PosL is  \SY{associative stacking} using the structure induced by tuple concatenation.
\end{lemma}
\begin{proof}
    Analogously to what we did for \SetL, we can define a multiplication operation $\mtimescatob$ in \PosL.
    \todotext{JL: the definition below is incorrect -- the objects in our category are not posets!}

    Given two objects $\posA = \cObj{\posA_1,\ldots, \posA_n}$ and $\posB = \cObj{\posB_1,\ldots, \posB_n}$, we define
    \begin{equation}
        \label{eq:multposstack}
        % \begin{aligned}
        \cObj{\posA_1,\ldots, \posA_n} \mtimescatob \cObj{\posB_1,\ldots, \posB_n} \definedas
        % \tup{\cObj{\posAsetn{1},\ldots, \posAsetn{m}} \cprod \cObj{\posBsetn{1},\ldots, \posBsetn{n}}, \posleqof{\cObj{\posA_1,\ldots, \posA_m}\cprod \cObj{\posB_1,\ldots, \posB_n}}} \\
        \cObj{\posA_1,\ldots, \posA_n, \posB_1,\ldots, \posB_n}.
    \end{equation}
    %    where the order is defined as
    %    \begin{equation}\label{eq:PosL-order}
    %        \prfdoubleperiod{
    %            \pars{\maketupconcat{\posAnel{1}, \posBnel{1}}}
    %            \ %
    %            \posleqof{\posA\mtimescatob \posB}
    %            \ %
    %            \pars{\maketupconcat{\posAnel{2}, \posBnel{2}}}
    %        }{
    %            \pars{\posAnel{1}\posleqof\posA\posAnel{2}}
    %            \ \booland\ %
    %            \pars{\posBnel{1}\posleqof\posB\posBnel{2}}
    %        }
    %    \end{equation}
    For the multiplication on morphisms, we define
    \begin{equation}
        \prfperiod{
            \mora\colon \posA \mto \posB
        }{
            \quad
        }{
            \mora\colon \posC \mto \posD
        }{
            \defmap{
                \pars{
                    \mora
                    \mtimescatmor
                    \morb
                }
            }{
                \posA\mtimescatob\posC
            }{
                \mto
            }{
                \posB\mtimescatob\posD
            }{
                \maketupconcat{\posAel, \posCel}
            }{
                \maketupconcat{\mora(\posAel), \morb(\posCel)}
            }
        }
    \end{equation}
    We need to check that the expression $\maketupconcat{\mora(\posAel), \morb(\posCel)}$ is monotone.
    This can be easily seen because the order on $\posB\mtimescatob\posD$ is akin to a product order.
    The proof for associativity of $\mtimescatmor$ is the same as in the proof of $\SetL$ (\cref{lem:SetL-is-associative-stacking}).
\end{proof}
% ~$\tup{\posAnel{1},\ldots, \posAnel{m},\posBnel{1},\ldots, \posBnel{n}} \posleqof{\cObj{\posA_1,\ldots, \posA_m}\mtimescatob \cObj{\posB_1,\ldots, \posB_n}} \tup{\posAnel{1}',\ldots, \posAnel{m}',\posBnel{1}',\ldots, \posBnel{n}'}$ if and only if~$\posAnel{i}\posleqof{\posA_i}\posAnel{i}'$ and~$\posBnel{j}\posleqof{\posB_j}\posBnel{j}'$ for all~$i\setin \makeset{1,\ldots, m}$,~$j\setin \makeset{1,\ldots, n}$.
% % Note that this construction is analogous to the product poset.
% The resulting \SY{poset} of tuples is written as~$\cObj{\posA_1,\ldots, \posA_m,\posB_1,\ldots, \posB_n}$.
% We now want to show that one can define a structure for which~\PosL is  \SY{associative stacking}.
% To do so, we define the stacking operation on objects as the multiplication operation~$\cprod$ defined in \cref{eq:multposstack}.

% The stacking operation on morphisms is defined as
% \begin{equation}
%     \defmapcomma{
%         \mtimescatmor
%     }{
%         \Mor_{\PosL} \cartprod \Mor_{\PosL}
%     }{
%         \mto
%     }{
%         \Mor_{\PosL}
%     }{
%         \tup{\mora,\morb}
%     }{
%         \mora \cprod \morb
%     }
% \end{equation}
% where~$(\mora \mtimescatmor \morb)(\tup{\posAnel{1},\ldots, \posAnel{m}} \tupconcat \tup{\posBnel{1},\ldots, \posBnel{n}})\definedas \mora(\tup{\posAnel{1},\ldots, \posAnel{m}})\tupconcat \morb(\tup{\posBnel{1},\ldots, \posBnel{n}})$.

% \begin{lemma}[\PosL is associative stacking]\label{lem:PosL-associative-stacking}
%     \PosL, equipped with the aforementioned stacking operations on objects and morphisms, is  \SY{associative stacking}.
% \end{lemma}

% \begin{proof}
%     % We show that~\PosL is an  \SY{associative stacking} semicategory.
%     We first show that it is a stacking category, by showing that the stacking operations fulfill \cref{def:simple-stacking-semi-cat}, and then showing the compatibility condition.
%     The stacking operation on objects fulfills \cref{def:simple-stacking-semi-cat}, since, given two \SY{posets} of tuples it returns a \SY{poset} of tuples.
%     To show that the stacking operation on morphism fulfills \cref{def:simple-stacking-semi-cat}, we need to show that given monotone~$\mora,\morb$, the map~$\mapa \mtimescatmor \mapb$ is monotone.
%     This is easy to show, by leveraging the monotonicity of~$\mora$ and~$\morb$.
%     For brevity, in the following we write~$\posAnel{}=\tup{\posAnel{1},\ldots, \posAnel{m}}$,~$\posAnel{}'=\tup{\posAnel{1}',\ldots, \posAnel{m}'}$,~$\posBnel{}=\tup{\posBnel{1},\ldots, \posBnel{n}}$, and~$\posBnel{}'=\tup{\posBnel{1}',\ldots, \posBnel{n}'}$.
%     We have
%     \begin{equation}
%         \prfcomma{\posAnel{}\tupconcat \posBnel{} \posleqof{\cObj{\posA_1,\ldots, \posA_m,\posB_1,\ldots, \posB_n}} \posAnel{}'\tupconcat \posBnel{}'}
%         {\prftree{
%                 \left(\posAnel{} \posleqof{\cObj{\posA_1,\ldots, \posA_m}} \posAnel{}'\right) \booland \left(\posBnel{} \posleqof{\cObj{\posB_1,\ldots, \posB_n}} \posBnel{}'\right)}
%             {
%                 \prftree{
%                     \left(\mora(\posAnel{}) \posleqof{\cObj{\posC_1,\ldots, \posC_o}} \mora(\posAnel{}')\right) \booland \left(\morb(\posBnel{})\posleqof{\cObj{\posD_1,\ldots, \posD_t}} \morb(\posBnel{}')\right)}
%                 {
%                     \prftree{
%                         {\mora(\posAnel{}) \tupconcat \morb(\posBnel{})  \posleqof{\cObj{\posC_1,\ldots, \posC_o, \posD_1,\ldots, \posD_t}} \mora(\posAnel{}') \tupconcat \morb(\posBnel{}')}}
%                     {(\mora \mtimescatmor \morb)(\posAnel{}\tupconcat \posBnel{}) \posleqof{\cObj{\posC_1,\ldots, \posC_o,\posD_1,\ldots, \posD_t}} (\mora \mtimescatmor \morb)(\posAnel{}'\tupconcat \posBnel{}')
%                     }}}
%         }
%     \end{equation}
%     which shows monotonicity.
%     Furthermore, the two stacking operations are compatible:
%     \begin{equation}
%         \prfperiod{\mora\colon \cObj{\posA_1,\ldots,\posA_m}\mto \cObj{\posC_1,\ldots,\posC_o}}
%         {\morb\colon \cObj{\posB_1,\ldots,\posB_n}\mto \cObj{\posD_1,\ldots,\posD_t}}
%         {\mora \mtimescatmor \morb \colon \cObj{\posA_1, \ldots, \posA_m, \posB_1, \ldots, \posB_n} \mto \cObj{\posC_1, \ldots, \posC_o, \posD_1, \ldots, \posD_t}}
%     \end{equation}
%     %
%     The proof of associativity is parallel to the proof for~\SetL.
% \end{proof}
%



\subsection{The category~$\DPL$}

Analogously, we define the category $\DPL$.

\begin{definition}[The category \DPL]\label{def:DPL}
    \SYNDEF{category of tuple-posets and design problems}
    The category~$\DPL$ is the \SY{subcategory} of \DP where the objects are \SY{posets} of tuple posets (objects of $\PosL$).
    % \begin{enumerate}
    %     \item \emph{Objects:}
    %           \SY{posets} of tuples.
    %     \item \emph{Morphisms:}
    %           A morphism from~$\posA =\cObj{\posA_1,\ldots, \posA_m}$ to ~$\posB = \cObj{\posB_1,\ldots, \posB_m}$ is design problem:
    %           \begin{equation}
    %               \adp\colon \posA\op\cartprod \posB \toinPos \Bool.
    %           \end{equation}
    %     \item \emph{Composition:}
    %           Composition is the usual composition of \SY{design problems}.
    %     \item \emph{Identities:}
    %           The identity morphism on an object~$\posA=\cObj{\posA_1,\ldots,\posA_n}$ is given by the identity \SY{design problem}~$\catid_\posA$.
    % \end{enumerate}
\end{definition}
% \begin{remark}
%     In other words, \DPL is the subcategory of \DP in which objects are of the form presented in \cref{eq:posetoftuples}.
% \end{remark}
\begin{lemma}
    \DPL is \SY{associative stacking} using the structure induced by tuple concatenation.
\end{lemma}
\begin{widepar}
    \begin{proof}
        For the stacking operation $\mtimescatob$ on objects, we use $\mtimescatobof{\DPL}\definedas\mtimescatobof{\PosL}$.
        For stacking morphisms, we define $\mtimescatmor$ by
        % Given two \SY{design problems} % \begin{equation}
        %     \begin{aligned}
        %         \adpa\colon \posA\posop \Ptimes \posC & \toinPos \Bool, \\
        %         \adpb\colon \posB\posop \Ptimes \posD & \toinPos \Bool,
        %     \end{aligned}
        % \end{equation}
        % we define
        \begin{equation}
            \label{eq:dpmulti}
            \prfperiod{
                \adpa\colon \posAop\Ptimes \posC \mto \Bool
            }{\qquad}{
                \adpb\colon \posB\posop\Ptimes \posD \mto \Bool,
            }{
                \defmap{
                    \adpa \mtimescatmor \adpb
                }{
                    \pars{\posA \mtimescatob \posB}\op
                    \Ptimes
                    \pars{\posC \mtimescatob \posD}
                }{
                    \mto
                }{
                    \Bool
                }{
                    \tup{\maketupconcat{a^*, c^*},\ \maketupconcat{b, d}}
                }{
                    \adpa(a^*,c) \booland \adpb(c^*, d)
                }
            }
        \end{equation}
        Note that this is a valid definition of a \SY{design problem} because the expression $\adpa(a^*,c) \booland \adpb(c^*, d)$ is monotone, as may readily be checked.

        To show associativity, consider three DPs
        \begin{equation}
            \adpa\colon \posAop\Ptimes \posC \toinPos \Bool, \qquad
            \adpb\colon \posB\posop\Ptimes \posD \toinPos \Bool, \qquad
            \adpc\colon \posE\posop\Ptimes \posF \toinPos \Bool.
        \end{equation}
        We compute $(\adpa \mtimescatmor \adpb) \mtimescatmor \adpc$
        and $\adpa \mtimescatmor (\adpb \mtimescatmor \adpc)$ according to the recipe in~\cref{eq:dpmulti}:
        \begin{equation}
            \defmapcomma{
                \parslight{\adpa \mtimescatmor \adpb} \mtimescatmor \adpc
            }{
                \pars{\parslight{\posA \mtimescatob \posB} \mtimescatob \posE}\op
                \Ptimes
                \pars{\parslight{\posC \mtimescatob \posD} \mtimescatob \posF}
            }{
                \toinPos
            }{
                \Bool
            }{
                \tup{
                    \parslight{\maketupconcat{a^*, c^*}} \tupconcat e^*
                    ,\
                    \parslight{\maketupconcat{b, d}} \tupconcat f
                }
            }{
                \parslight{
                    \adpa(a^*,c) \booland \adpb(c^*, d)
                }
                \ \booland\
                \adpc(e^*, f)
            }
        \end{equation}
        \begin{equation}
            \defmapperiod{
                \adpa \mtimescatmor \parslight{\adpb \mtimescatmor \adpc}
            }{
                \pars{
                    \posA \mtimescatob
                    \parslight{
                        \posB \mtimescatob \posE
                    }
                }\op
                \Ptimes
                \pars{
                    \posC \mtimescatob
                    \parslight{
                        \posD \mtimescatob \posF
                    }
                }
            }{
                \toinPos
            }{
                \Bool
            }{
                \tup{
                    \maketupconcat{a^*, \parslight{\maketupconcat{c^*, e^*})}},\
                    \maketupconcat{b, \parslight{\maketupconcat{d, f}}}
                }
            }{
                \adpa(a^*,c)
                \ \booland\
                \parslight{
                    \adpb(c^*, d)
                    \booland
                    \adpc(e^*, f)
                }
            }
        \end{equation}
        Because the operations $\mtimescatob$ and $\booland$ are \SY{associative}, we can erase all the light parentheses in the formulas, and we find that $(\adpa \mtimescatmor \adpb) \mtimescatmor \adpc$ and
        $\adpa \mtimescatmor (\adpb \mtimescatmor \adpc)$ are equal to the design problem
        \begin{equation}
            \defmapperiod{
                \adpa \mtimescatmor \adpb \mtimescatmor \adpc
            }{
                \pars{
                    \posA \mtimescatob
                    \posB \mtimescatob \posE
                }\op
                \Ptimes
                \pars{
                    \posC \mtimescatob
                    \posD \mtimescatob \posF
                }
            }{
                \toinPos
            }{
                \Bool
            }{
                \tup{
                    \maketupconcat{a^*, c^*, e^*},\
                    \maketupconcat{b, d, f}
                }
            }{
                \adpa(a^*,c)
                \booland
                \adpb(c^*, d)
                \booland
                \adpc(e^*, f)
            }
        \end{equation}
    \end{proof}
\end{widepar}

% where
% \begin{equation}
%     (\adpa \cproddp \adpb)  (\tup{\posAnel{}\opel,\posBnel{}\opel,\posCnel{},\posDnel{}})
%     \definedas\adpa(\tup{\posAnel{}\opel,\posCnel{}})\booland
%     \adpb(\tup{\posBnel{}\opel,\posDnel{}}).
% \end{equation}

% We now consider the stacking operation on objects we considered for~\PosL, and the stacking operation on morphism defined as
% \begin{equation}
%     \defmapcomma{
%         \mtimescatmor
%     }{
%         \Mor_{\DPL} \cartprod \Mor_{\DPL}
%     }{
%         \mto
%     }{
%         \Mor_{\DPL}
%     }{
%         \tup{\adpa,\adpb}
%     }{
%         \adpa \cproddp \adpb
%     },
% \end{equation}

% \begin{lemma}\label{lem:DPL-associative-stacking}
%     \DPL, equipped with the aforementioned stacking operations on objects and morphisms, is  \SY{associative stacking}.
% \end{lemma}

% \begin{proof}
%     We show that~$\DPL$ is an  \SY{associative stacking} semicategory.
%     We first show that it is a stacking category, by first showing that the stacking operations indeed fulfill \cref{def:simple-stacking-semi-cat}, and then showing the compatibility condition.
%     The stacking operation on objects fulfills \cref{def:simple-stacking-semi-cat}, since, given two \SY{posets} of tuples it returns a \SY{poset} of tuples.
%     To show that the stacking operation on morphism fulfills \cref{def:simple-stacking-semi-cat}, we need to show that given \SY{design problems} ~$\adpa,\adpb$, the morphism~$\adpa \mtimescatmor \adpb$ is a design problem.
%     This is easy to show, by leveraging the fact that~$\adpa$ and~$\adpb$ are \SY{design problems}.
%     We know that
%     \begin{equation}
%         \prfcomma{\tup{\posAnel{}\opel,\posBnel{}\opel,\posCnel{},\posDnel{}} \posleqof{\posA\posop \Ptimes \posB\posop \Ptimes \posC\Ptimes \posD} \tup{{\posAnel{}'}\opel,{\posBnel{}'}\opel,\posCnel{}',\posDnel{}'}}
%         {(\posAnel{}\opel \posleqof{\posA\op} {\posAnel{}'}\opel) \booland (\posBnel{}\opel \posleqof{\posB\op} {\posBnel{}'}\opel)\booland (\posCnel{}\posleqof{\posC}\posCnel{}')\booland (\posDnel{}\posleqof{\posD}\posDnel{}')}
%     \end{equation}
%     and that~$\booland$ is a monotone operator.
%     Therefore,~$\adpa \cproddp \adpb$ is a design problem.
%     Furthermore, the two stacking operations are compatible:
%     \begin{equation}
%         \prfperiod{\adpa\colon \posA\posop \Ptimes \posC \toinPos \Bool}
%         {\adpb\colon \posB\posop \Ptimes \posD\toinPos \Bool}
%         {\adpa \cproddp \adpb \colon \posA\posop \Ptimes \posB\posop \Ptimes\posC\Ptimes\posD\toinPos \Bool}
%     \end{equation}

%     \begin{comment}
%     For brevity, we write~$\posAnel{}=\tup{\posAnel{1},\ldots, \posAnel{m}}$,~$\posBnel{}=\tup{\posBnel{1},\ldots, \posBnel{n}}$,~$\posCnel{}=\tup{\posCnel{1},\ldots, \posCnel{o}}$, and~$\posDnel{}=\tup{\posDnel{1},\ldots, \posDnel{t}}$.
%     We have:
%     \begin{equation}
%         \prfperiod{
%         \posAnel{}\tupconcat \posBnel{}\posleqof{\cObj{\posA_1, \ldots, \posA_m, \posB_1, \ldots, \posB_n}} \posAnel{}'\tupconcat \posBnel{}'
%         }
%         {\prftree{(\posAnel{} \posleqof{\cObj{\posA_1,\ldots, \posA_m}} \posAnel{}') \booland (\posBnel{} \posleqof{\cObj{\posB_1,\ldots, \posB_n}} \posBnel{}')}
%         {\prftree{(\adpa(\posAnel{}\opel,\posCnel{})\posgeqof{\Bool} \adpa({\posAnel{}'}\opel,\posCnel{}))
%                 \booland
%                 (\adpb(\posBnel{}\opel,\posDnel{})\posgeqof{\Bool} \adpb({\posBnel{}'}\opel,\posDnel{}))}
%             {
%                 \prftree{
%                     (\adpa({\posAnel{}}\opel,\posCnel{})\booland \adpb(\posBnel{}\opel,\posDnel{}))
%                     \posgeqof{\Bool}
%                     (\adpa({\posAnel{}'}\opel,\posCnel{})\booland \adpb({\posBnel{}'}\opel,\posDnel{}))
%                 }
%                 {(\adpa \cprod \adpb)((\posAnel{}\tupconcat \posBnel{})\opel,\posCnel{}\tupconcat \posDnel{})\posgeqof{\Bool}
%                     (\adpa \cprod \adpb)((\posAnel{}'\tupconcat \posBnel{}')\opel,\posCnel{}\tupconcat \posDnel{})} }
%             }
%         }
%     \end{equation}
%     For the second iteration, we have:
%     \begin{equation}
%         \prfperiod{
%         \posCnel{}\tupconcat \posDnel{}\posleqof{\cObj{\posC_1, \ldots, \posC_o, \posD_1, \ldots, \posD_t}} \posCnel{}'\tupconcat \posDnel{}'
%         }
%         {\prftree{(\posCnel{} \posleqof{\cObj{\posC_1,\ldots,\posC_o}} \posCnel{}') \booland (\posDnel{}\posleqof{\cObj{\posD_1,\ldots, \posD_t}} \posDnel{}')}
%         {\prftree{(\adpa(\posAnel{}\opel,\posCnel{})\posleqof{\Bool}
%                 \adpa(\posAnel{}\opel,\posCnel{}'))
%                 \booland
%                 (\adpb(\posBnel{}\opel,\posDnel{})\posleqof{\Bool}
%                 \adpb(\posBnel{}\opel,\posDnel{}'))}
%             {
%                 \prftree{
%                     (\adpa(\posAnel{}\opel,\posCnel{})\booland \adpb(\posBnel{}\opel,\posDnel{}))
%                     \posleqof{\Bool}
%                     (\adpa(\posAnel{}\opel,\posCnel{}')\booland \adpb(\posBnel{}\opel,\posDnel{}'))
%                 }
%                 {(\adpa \cprod \adpb)((\posAnel{}\tupconcat \posBnel{})\opel,\posCnel{}\tupconcat \posDnel{})
%                     \posleqof{\Bool}
%                     (\adpa \cprod \adpb)((\posAnel{}\tupconcat \posBnel{})\opel,\posCnel{}'\tupconcat \posDnel{}')} }
%             }
%         }
%     \end{equation}
%     Therefore,~$\adpa \cprod \adpb$ is really a design problem.
%     \end{comment}
%     Associativity can be shown following what we did for~$\cObj{\Set}$.
% \end{proof}

\begin{example}[\SetL is a functorial stacking semicategory]
    \label{ex:setfunstack}
    We want to show that~\SetL is a \SY{functorial stacking semicategory}.

    Consider four morphisms
    \begin{equation}
        \begin{aligned}
            \mora\colon \setA \mto \setB, & \quad \morc\colon \setB \mto \setC, \\
            \morb\colon \setD \mto \setE, & \quad \mord\colon \setE \mto \setF.
            \\
        \end{aligned}
    \end{equation}
    We want to show that
    \begin{equation}
        (\mora \mtimescatmor \morb )
        \mthen (\morc \mtimescatmor \mord )
        =
        (\mora \mthen \morc ) \mtimescatmor (\morb \mthen \mord ),
    \end{equation}
    We show this by showing that, for any $\setAel\tupconcat \setDel\setin \setA\cprod\setD$:
    \begin{equation}
        \begin{aligned}
            ((\mora \mtimescatmor \morb )\mthen (\morc \mtimescatmor \mord ))(\setAeln{}\tupconcat \setDeln{})
             & = (\morc \mtimescatmor \mord)(\mora(\setAeln{})\tupconcat \morb(\setDeln{})) \\
             & =\morc(\mora(\setAeln{}))\tupconcat \mord(\morb(\setDeln{})) \\
             & =(\mora \mthen \morc)(\setAeln{})\tupconcat (\morb\mthen \mord)(\setDeln{}) \\
             & =((\mora \mthen \morc ) \mtimescatmor (\morb \mthen \mord ))(\setAeln{}\tupconcat \setDeln{}).
        \end{aligned}
    \end{equation}
\end{example}

\begin{exercise}
    \label{ex:the-usual-suspects-are-functorial-stacking-categories}
    Prove that \PosL is a functorial stacking category.
    Hint: You have to define the functor $\mtimescat$ and check that it satisfies the two equations in \cref{def:functorial-stacking-cat}.
\end{exercise}
\begin{solution}
    \todotext{Solution for \cref{ex:the-usual-suspects-are-functorial-stacking-categories}}
\end{solution}

\begin{gradedexercise}[\exname{RelFunStack}]
    \label{ex:RelFunStack}
    Prove that the structure defined in \cref{ex:relassocstack} makes~\RelL a \SY{functorial stacking semicategory}.
\end{gradedexercise}

\solutionof{RelFunStack}

\begin{lemma}[\DPL is a functorial stacking semicategory]\label{lem:DPL-functorial-stacking}
    \DPL, equipped with the aforementioned stacking operations on objects and morphisms, is \SY{functorial stacking semicategory}.
\end{lemma}

\begin{proof}
    Consider
    \begin{equation}
        \begin{aligned}
            \adpa\colon \posAop\Ptimes \posC     & \toinPos \Bool \\
            \adpb\colon \posB\posop\Ptimes \posD & \toinPos \Bool \\
            \adpc\colon \posC\posop\Ptimes \posE & \toinPos \Bool \\
            \adpd\colon \posD\posop\Ptimes \posF & \toinPos \Bool
        \end{aligned}
    \end{equation}
    We want to prove that
    \begin{equation}
        (\adpa \dpthen \adpc)
        \mtimescatmor (\adpb \dpthen \adpd)=(\adpa \mtimescatmor \adpb) \dpthen (\adpc \mtimescatmor \adpd).
    \end{equation}
    %For brevity, in the following we write
    %\begin{align*}
    %    \posAnel{} & =\tup{\posAnel{1},\ldots, \posAnel{i}}, &
    %    \posBnel{} & =\tup{\posBnel{1},\ldots, \posBnel{l}} \\
    %    \posCnel{} & =\tup{\posCnel{1},\ldots, \posCnel{m}}, &
    %    \posDnel{} & =\tup{\posDnel{1},\ldots, \posDnel{n}} \\
    %    \posEnel{} & =\tup{\posEnel{1},\ldots, \posEnel{o}}, &
    %    \posFnel{} & =\tup{\posFnel{1},\ldots, \posFnel{v}}
    %\end{align*}
    We start from the left-hand side.
    We have
    \begin{equation}
        (\adpa \dpthen \adpc)(\posAnel{}\opel,\posEnel{})
        =\bigvee_{\posCnel{}\setin \posC}
        \adpa(\posAnel{}\opel, \posCnel{}) \booland \adpc(\posCnel{}\opel,\posEnel{})
    \end{equation}
    and
    \begin{equation}
        \begin{aligned}
            (\adpb \dpthen \adpd)(\posBnel{}\opel,\posFnel{})
            =\bigvee_{\posDnel{}\setin \posD}
            \adpb(\posBnel{}\opel, \posDnel{}) \booland \adpd(\posDnel{}\opel\posFnel{})
        \end{aligned}
    \end{equation}
    Therefore, we know
    \begin{equation}
        \begin{aligned}
             & ((\adpa \dpthen \adpc)\mtimescatmor (\adpb \dpthen \adpd))
            ((\posAnel{}\tupconcat \posBnel{})\opel,\posEnel{}\tupconcat \posFnel{}) \\
             & =\bigvee_{\posCnel{}\setin \posC}
            \adpa(\posAnel{}\opel, \posCnel{}) \booland \adpc(\posCnel{}\opel,\posEnel{}) \booland
            \bigvee_{\posDnel{}\setin \posD} \adpb(\posBnel{}\opel, \posDnel{}) \booland \adpd(\posDnel{}\opel,\posFnel{}).
        \end{aligned}
    \end{equation}

    On the other hand, we have
    \begin{equation}
        (\adpa \mtimescatmor \adpb)
        ((\posAnel{}\tupconcat \posBnel{})\opel, \posCnel{}\tupconcat \posDnel{})
        =\adpa(\posAnel{}\opel,\posCnel{}) \booland \adpb(\posBnel{}\opel,\posDnel{})
    \end{equation}
    and
    \begin{equation}
        (\adpc \mtimescatmor \adpd)
        ((\posCnel{}\tupconcat \posDnel{})\opel, \posEnel{}\tupconcat \posFnel{})
        =\adpc(\posCnel{}\opel,\posEnel{}) \booland \adpd(\posDnel{}\opel,\posFnel{})
    \end{equation}
    Therefore, we know
    \begin{equation}
        \begin{aligned}
             & ((\adpa \mtimescatmor \adpb)\dpthen (\adpc \mtimescatmor \adpd))((\posAnel{}\tupconcat \posBnel{})\opel,\posEnel{}\tupconcat \posFnel{}) \\
             & =\bigvee_{\posCnel{}\tupconcat \posDnel{}\setin \cObj{\posC,\posD}}
            (\adpa \mtimescatmor \adpb)((\posAnel{}\tupconcat \posBnel{})\opel, \posCnel{}\tupconcat \posDnel{})
            \booland (\adpc \mtimescatmor \adpd)((\posCnel{}\tupconcat \posDnel{})\opel, \posEnel{}\tupconcat \posFnel{}) \\
             & =\bigvee_{\posCnel{}\tupconcat \posDnel{}\setin \cObj{\posC,\posD}}
            \adpa(\posAnel{}\opel,\posCnel{}) \booland \adpb(\posBnel{}\opel,\posDnel{})
            \booland
            \adpc(\posCnel{}\opel,\posEnel{}) \booland \adpd(\posDnel{}\opel,\posFnel{}) \\
             & =\bigvee_{\posCnel{}\setin \posC}
            \adpa(\posAnel{}\opel, \posCnel{}) \booland \adpc(\posCnel{}\opel,\posEnel{}) \booland
            \bigvee_{\posDnel{}\setin \posD} \adpb(\posBnel{}\opel, \posDnel{}) \booland \adpd(\posDnel{}\opel,\posFnel{}),
        \end{aligned}
    \end{equation}
    proving the statement for any \SY{posets}~$\posA,\posB,\posC,\posD,\posE,\posF$ (and hence, also for \SY{posets} of tuples).

\end{proof}


